\section{Astronomical Null Models}\label{sec:astro-null}

A recurrent numeral becomes scientifically interesting only after comparison
with simple physical processes that could generate it without cultural
transmission or hidden formal coding.  The present section therefore treats
astronomy as a source of \emph{null models}, not as a post-hoc confirmation
engine.

\subsection{Three-day darkness as a lunar visibility null}

Near conjunction the Moon lies close to the Sun and is invisible to the naked
eye.  Modern visibility work gives a typical invisibility interval spanning
roughly $2.5$--$4.5$ days, depending on geometry and observing conditions.
Consequently,
\[
3\ \mathrm{days}\in[2.5,4.5]\ \mathrm{days}.
\]
This matters because a mythic pattern
\[
\mathrm{visible/living}\to\mathrm{dark/liminal}\to
\mathrm{return}\quad\text{after roughly three days}
\]
can arise from a very ordinary astronomical recurrence.  A three-day marker
therefore has low discriminatory power unless the surrounding event graph is
also preserved.

This null model does not imply that any particular resurrection or underworld
myth is lunar.  It states the opposite methodological point: a physical process
already generates a similar timescale, so stronger evidence is required before
claiming a special cross-cultural relation.

\subsection{Seven days and a quarter lunation}

Using the mean synodic month
\[
T_{\mathrm{syn}}=29.530588\ \mathrm{d},
\]
one quarter is
\[
Q_{\mathrm{Moon}}=\frac{T_{\mathrm{syn}}}{4}
=7.382647\ \mathrm{d}.
\]
Thus seven days is numerically close to a lunar quarter, with an offset
\[
|Q_{\mathrm{Moon}}-7\ \mathrm{d}|=0.382647\ \mathrm{d}.
\]
This is a useful astronomical comparison, but it is deliberately not promoted
to an origin claim.  Modern historical scholarship has noted proposed links
between the Israelite seven-day week and Babylonian calendrical phenomena and
has also judged simple derivations unconvincing.  Therefore the repository
records this only as \emph{numerical proximity plus a plausible astronomical
scale}, not as historical causation.

\subsection{A positive control: Egyptian lunar mythography}

A stronger kind of evidence exists in ancient Egypt.  Coffin Texts spells
154--160 have been analysed as a chronologically ordered account of the lunar
month expressed through religious and mythological language: lunar
invisibility, waxing, events around full moon, waning, the last crescent, and
renewed conjunction.  This serves as a positive control for the programme:
ancient religious texts \emph{can} encode ordered astronomical phenomena.

The logical implication is only
\[
\exists\,\mathcal M:\quad
\mathcal M_{\mathrm{mythic}}\leftrightarrow
\mathcal C_{\mathrm{astronomical}}.
\]
It does not imply
\[
\forall\,\mathcal M:\quad
\mathcal M_{\mathrm{mythic}}=\mathcal C_{\mathrm{astronomical}}.
\]

\subsection{Mesopotamian quantitative capability}

Mesopotamian astronomy provides a second positive control at the level of
method.  Ancient scholars numerically schematized periodic lunar visibility,
daylight variation, eclipses, synodic appearances of planets, and first lunar
visibility.  Therefore the historical cultures in which several relevant
myths circulated possessed genuine quantitative concepts of celestial
periodicity.  This supports the \emph{capability} premise of The Mathematics
of Myth while leaving every specific myth--astronomy identification open to
source-level falsification.

\subsection{Consequence for the transition hypothesis}

The current hierarchy of evidence is therefore
\[
\begin{aligned}
\text{same number}
&<\text{astronomically natural timescale}\\
&<\text{same ordered transition graph}\\
&<\text{source-attested astronomical mapping}.
\end{aligned}
\]
A robust claim must move upward through this hierarchy.  In particular, the
$3$-threshold and $6\to7$ hypotheses are not strengthened merely because lunar
astronomy contains nearby numbers; rather, lunar astronomy provides a null
against which the structural claims must outperform coincidence.

